\phantomsection\label{sec:research-plan}
\section*{Research Plan}
\subsection*{Established Research: Spatiotemporal Data Mining}
My established research is in spatiotemporal data mining, the extraction and use of patterns from large-scale data that carries spatial information and changes over time, usually sourced from transportation scenarios.
A theme across my works is representation learning, which learns task-agnostic representations of data so that a separate model is not needed for each task.
Representation learning is a promising path to foundational models for spatiotemporal data mining that generalize across data sources and tasks, as demonstrated in my works, including UniTE, UVTM, and TransferTraj.

\subsection*{Emerging Direction: AI for Materials Science}
Building on my established research, I am developing an interdisciplinary direction that applies data mining to materials science, with the goal of advancing AI for materials science.
This direction is a natural progression of my established research, because materials, such as glasses, are themselves spatial data with a dynamic structure, and mining them needs the same representation learning I develop for spatiotemporal data.
This direction is already underway through my Villum-funded postdoctoral project on \textit{inverse design of materials using diffusion probabilistic models}, in collaboration with Morten M. Smedskjaer's group, one of the leading groups in disordered materials science.

\vspace{\entryskip}
The promise of this direction is that AI can discover new materials with tailored properties in a vast design space, which traditional trial-and-error approaches are too slow and costly to explore.
My work in this direction has already produced AMDEN, a diffusion model for the inverse design of amorphous materials, published with me as joint first author in Advanced Materials in 2026.
I plan to deepen and expand this direction and to become a leading researcher in AI for materials science, with a focus on disordered materials, an area with strong research and practical potential that remains understudied.

\subsection*{Broader Impact}
Disordered materials, such as glasses, underpin a broad range of technologies, including energy storage, energy production, and advanced optics.
AI can replace many resource-heavy laboratory trial-and-error experiments and speed up the discovery of new disordered materials with tailored properties.
Faster discovery of these materials supports the green transition goals of the EU and the UN, including the Sustainable Development Goals on responsible production and climate action.
The disordered materials group I collaborate with can prepare and validate promising candidates in the laboratory, so successful designs can move toward real use.

\vspace{\entryskip}
By connecting computer science and materials science, this direction also builds lasting interdisciplinary research capacity and a collaboration between the two fields that benefits both in the long term.
